\documentclass[12pt, a4paper]{article}

\usepackage[margin=2.5cm, headheight=14.5pt]{geometry}
\usepackage{amsmath, amssymb}
\usepackage{hyperref}
\usepackage{xcolor}
\usepackage{listings}
\usepackage{booktabs}
\usepackage{enumitem}
\usepackage{tcolorbox}
\usepackage{fancyhdr}
\usepackage{titlesec}
\usepackage{microtype}
\usepackage{tikz}
\usetikzlibrary{shapes.geometric, arrows.meta, positioning, fit}

\tcbuselibrary{skins, breakable, listings}

% ── Colours ──────────────────────────────────────────────────────────────────
\definecolor{codeblue}{RGB}{0, 70, 140}
\definecolor{codebg}{RGB}{246, 248, 250}
\definecolor{borderblue}{RGB}{180, 210, 240}
\definecolor{funcgreen}{RGB}{10, 100, 50}
\definecolor{noteamber}{RGB}{180, 100, 0}
\definecolor{lightamber}{RGB}{255, 248, 220}
\definecolor{lightgreen}{RGB}{230, 255, 240}
\definecolor{lightblue}{RGB}{230, 242, 255}

% ── Hyperref ─────────────────────────────────────────────────────────────────
\hypersetup{colorlinks=true, linkcolor=codeblue, urlcolor=funcgreen,
  pdftitle={C++ Code Walkthrough --- Linear Algebra Analyzer}}

% ── Section formatting ───────────────────────────────────────────────────────
\titleformat{\section}{\large\bfseries\color{codeblue}}{}{0em}{\thesection\quad}
\titleformat{\subsection}{\normalsize\bfseries\color{codeblue!80!black}}{}{0em}{\thesubsection\quad}
\titleformat{\subsubsection}{\normalsize\bfseries\color{black}}{}{0em}{\thesubsubsection\quad}

% ── C++ listings style ───────────────────────────────────────────────────────
\lstdefinestyle{cpp}{
  language=C++,
  basicstyle=\footnotesize\ttfamily,
  keywordstyle=\bfseries\color{codeblue},
  commentstyle=\itshape\color{gray!70},
  stringstyle=\color{funcgreen},
  numberstyle=\tiny\color{gray},
  numbers=left, numbersep=6pt,
  frame=single, rulecolor=\color{borderblue},
  backgroundcolor=\color{codebg},
  breaklines=true, tabsize=4,
  showstringspaces=false,
  morekeywords={vector, string, tuple, pair, complex, auto},
  captionpos=b
}
\lstset{style=cpp}

% ── Callout boxes ────────────────────────────────────────────────────────────
\newtcolorbox{notebox}[1]{
  colback=lightamber, colframe=noteamber,
  fonttitle=\bfseries\small\color{noteamber}, title=#1,
  arc=3pt, boxrule=0.7pt, left=4pt, right=4pt}

\newtcolorbox{flowbox}[1]{
  colback=lightblue, colframe=codeblue,
  fonttitle=\bfseries\small, title=#1,
  arc=3pt, boxrule=0.7pt}

\newtcolorbox{funcbox}[1]{
  colback=lightgreen, colframe=funcgreen,
  fonttitle=\bfseries\small, title=#1,
  arc=3pt, boxrule=0.7pt, left=4pt, right=4pt}

% ── Header / Footer ──────────────────────────────────────────────────────────
\pagestyle{fancy}
\fancyhf{}
\lhead{\small\color{codeblue}C++ Code Walkthrough --- Linear Algebra Analyzer}
\rhead{\small\color{gray}\thepage}
\renewcommand{\headrulewidth}{0.4pt}

% ── Math macros ──────────────────────────────────────────────────────────────
\newcommand{\code}[1]{\texttt{\small#1}}
\newcommand{\fn}[1]{\texttt{\small\color{funcgreen}#1}}
\newcommand{\kw}[1]{\texttt{\small\bfseries\color{codeblue}#1}}

% ─────────────────────────────────────────────────────────────────────────────

\begin{document}

% ── Title ────────────────────────────────────────────────────────────────────
\begin{titlepage}
\centering\vspace*{2cm}
{\huge\bfseries\color{codeblue}C++ Code Walkthrough}\\[0.5em]
{\large Linear Algebra Analyzer --- Full Source Explanation}\\[2em]
\rule{0.7\textwidth}{0.5pt}\\[1.5em]
{\normalsize
  \texttt{linalg.h} $\cdot$ \texttt{linalg.cpp} $\cdot$ \texttt{main.cpp}\\[0.3em]
  984 lines total $\cdot$ C++17 $\cdot$ Standard Library only
}\\[1.5em]
\rule{0.7\textwidth}{0.5pt}\\[2em]
{\large Sarvesh S}\\[0.4em]
{\small IIT Madras \quad|\quad \texttt{ed23b042@smail.iitm.ac.in}}\\[3em]
\begin{tcolorbox}[colback=lightblue, colframe=codeblue, width=0.75\textwidth, arc=5pt]
\centering\small
This document explains \emph{exactly how the code works} ---
every function, every loop, every design choice.
Theory is covered separately in \texttt{linear\_algebra\_doc.pdf}.
\end{tcolorbox}
\vfill{\small\color{gray}\today}
\end{titlepage}

\tableofcontents
\newpage

% ─────────────────────────────────────────────────────────────────────────────
\section{Project Architecture}
% ─────────────────────────────────────────────────────────────────────────────

\subsection{File layout and responsibilities}

\begin{center}
\small
\begin{tabular}{llr}
\toprule
File & Responsibility & Lines \\
\midrule
\texttt{linalg.h}   & Class declaration, result-type structs, all \code{inline} predicates & 129 \\
\texttt{linalg.cpp} & Every algorithm implementation & 637 \\
\texttt{main.cpp}   & Interactive CLI, formatted output & 218 \\
\bottomrule
\end{tabular}
\end{center}

\subsection{The \texttt{Matrix} class}

Everything lives in one class, \code{Matrix}.  Its public data members are:
\begin{lstlisting}
int m, n;                              // rows, columns
vector<vector<double>> A;              // A[i][j] = element at row i, col j
static constexpr double EPS = 1e-9;   // global zero threshold
\end{lstlisting}

\code{A} is a \emph{row-major} 2-D vector: \code{A[i]} is the $i$-th row as a
\code{vector<double>}.  Swapping two rows (\code{std::swap(A[i], A[j])}) is
therefore $O(1)$ — it just swaps two pointers, not all the data.  This is
exploited heavily in Gauss elimination and LUP.

\subsection{Why result structs live \emph{outside} the class}

The three result types — \code{RREFResult}, \code{EigenResult},
\code{JordanInfo} — each contain a \code{Matrix} member:

\begin{lstlisting}
struct RREFResult { Matrix R; vector<int> pivot_cols; int rank; };
\end{lstlisting}

In C++, a nested struct defined \emph{inside} a class body is parsed before
the class is complete, so the incomplete type \code{Matrix} cannot appear as a
value member.  The fix: forward-declare the structs before the class so their
names are visible inside the class for return types, then define them
\emph{after} the class is complete:

\begin{lstlisting}
struct RREFResult;   // forward declaration
class Matrix { ... RREFResult rref() const; ... };
struct RREFResult { Matrix R; ... };   // full definition, Matrix is complete
\end{lstlisting}

\subsection{C++17 structured bindings}

Throughout the code, return values are destructured with \code{auto}:
\begin{lstlisting}
auto [R, pivot_cols, rank_] = A.rref();   // RREFResult
auto [L, U, P]              = A.LUP();    // tuple<Matrix,Matrix,Matrix>
auto [Q, Rm]                = A.QR();     // pair<Matrix,Matrix>
auto [U_s, S_s, V_s]        = A.SVD();    // tuple<Matrix,Matrix,Matrix>
\end{lstlisting}

\code{RREFResult} is a C++ aggregate (public members, no constructor), and
C++17 aggregate structured bindings work on those exactly like \code{tuple}.

% ─────────────────────────────────────────────────────────────────────────────
\section{Constructors and Static Helpers}
% ─────────────────────────────────────────────────────────────────────────────

\begin{lstlisting}
Matrix::Matrix(int m_, int n_, double val)
    : m(m_), n(n_), A(m_, vector<double>(n_, val)) {}
\end{lstlisting}

The initialiser list builds \code{A} as \code{m\_} rows each containing
\code{n\_} copies of \code{val}.  Calling \code{Matrix(3,4,0.0)} gives a
$3\times4$ zero matrix.  Default \code{val=0.0} so \code{Matrix(3,4)} also
gives zeros.

\begin{lstlisting}
Matrix::Matrix(const vector<vector<double>>& data)
    : m(data.size()), n(data.empty() ? 0 : data[0].size()), A(data) {}
\end{lstlisting}

Allows direct construction from a 2-D \code{vector}, e.g.
\code{Matrix(\{\{1,2\},\{3,4\}\})}.

\medskip
\begin{lstlisting}
Matrix Matrix::eye(int n) {
    Matrix I(n, n, 0.0);
    for (int i = 0; i < n; i++) I.A[i][i] = 1.0;
    return I;
}
\end{lstlisting}

\fn{eye(n)} starts from a zero matrix, then sets the diagonal.
\code{Matrix::zeros(r,c)} just calls the value constructor.
\fn{fromInput} reads from \code{std::cin} row-by-row.

% ─────────────────────────────────────────────────────────────────────────────
\section{Basic Operations}
% ─────────────────────────────────────────────────────────────────────────────

\subsection{Arithmetic operators}

\begin{lstlisting}
Matrix Matrix::operator*(const Matrix& B) const {
    if (n != B.m) throw runtime_error("*: size mismatch");
    Matrix C(m, B.n, 0.0);
    for (int i = 0; i < m; i++)
        for (int k = 0; k < n; k++)
            for (int j = 0; j < B.n; j++)
                C.A[i][j] += A[i][k] * B.A[k][j];
    return C;
}
\end{lstlisting}

The loop order is \code{i-k-j} rather than the textbook \code{i-j-k}.
The inner loop \code{j} walks along a row of \code{B} and a row of \code{C}
with \code{A[i][k]} fixed.  This is better for CPU cache: \code{B.A[k]} and
\code{C.A[i]} are both contiguous in memory.

\subsection{Transpose}

\begin{lstlisting}
Matrix Matrix::T() const {
    Matrix C(n, m);
    for (int i = 0; i < m; i++)
        for (int j = 0; j < n; j++)
            C.A[j][i] = A[i][j];   // note the flip
    return C;
}
\end{lstlisting}

Returns a \emph{new} matrix of size $n \times m$.  Chaining is possible:
\code{A.T() * A} computes $A^TA$.

\subsection{print()}

\begin{lstlisting}
void Matrix::print(const string& label, int prec) const {
    if (!label.empty()) cout << label << ":\n";
    for (int i = 0; i < m; i++) {
        cout << "  [";
        for (int j = 0; j < n; j++) {
            double v = (abs(A[i][j]) < EPS) ? 0.0 : A[i][j];
            cout << setw(prec+6) << fixed << setprecision(prec) << v;
        }
        cout << " ]\n";
    }
}
\end{lstlisting}

The line \code{double v = (abs(...) < EPS) ? 0.0 : A[i][j]} cleans up
floating-point noise before printing --- values smaller than $10^{-9}$ are
shown as \code{0.0000} instead of \code{-0.0000} or \code{1.23e-14}.

% ─────────────────────────────────────────────────────────────────────────────
\section{RREF --- The Central Engine}
% ─────────────────────────────────────────────────────────────────────────────

\code{rref()} is the most-called function in the project.  It is used by
\fn{rank}, \fn{nullSpace}, \fn{columnSpace}, \fn{rowSpace}, \fn{inv}, and
\fn{jordanForm}.

\begin{lstlisting}
RREFResult Matrix::rref() const {
    Matrix R = *this;          // work on a copy; don't modify A
    vector<int> pivot_cols;
    int cur_row = 0;

    for (int col = 0; col < n && cur_row < m; col++) {
        // (1) Find pivot: largest |entry| in column, from cur_row down
        int piv = -1;
        double best = EPS;
        for (int r = cur_row; r < m; r++) {
            if (abs(R.A[r][col]) > best) { best = abs(R.A[r][col]); piv = r; }
        }
        if (piv == -1) continue;   // entire column is zero -> free variable

        // (2) Swap pivot row to cur_row
        swap(R.A[cur_row], R.A[piv]);

        // (3) Scale: make leading entry exactly 1
        double sc = R.A[cur_row][col];
        for (int j = 0; j < n; j++) R.A[cur_row][j] /= sc;

        // (4) Eliminate: zero out ALL other rows in this column
        for (int r = 0; r < m; r++) {
            if (r == cur_row) continue;
            double f = R.A[r][col];
            if (abs(f) < EPS) continue;
            for (int j = 0; j < n; j++) R.A[r][j] -= f * R.A[cur_row][j];
        }
        pivot_cols.push_back(col);
        cur_row++;
    }
    // (5) Snap near-zero to exact zero
    for (int i = 0; i < m; i++)
        for (int j = 0; j < n; j++)
            if (abs(R.A[i][j]) < EPS) R.A[i][j] = 0.0;

    return {R, pivot_cols, (int)pivot_cols.size()};
}
\end{lstlisting}

\begin{notebox}{Step-by-step trace for a $3\times4$ example}
Suppose $A = \begin{pmatrix}1&2&3&4\\2&4&6&8\\0&1&2&3\end{pmatrix}$:
\begin{itemize}[noitemsep]
  \item \textbf{col=0:} Largest in col 0 is row 1 (\code{|2|>|1|>|0|}).
        Swap rows 0$\leftrightarrow$1. Scale: divide row 0 by 2.
        Eliminate: row 1 $-= 1 \cdot$ row 0; row 2 unchanged.
        Record \code{pivot\_cols = \{0\}}, \code{cur\_row = 1}.
  \item \textbf{col=1:} Best below cur\_row=1 is row 2 ($|1| > |0|$).
        Swap rows 1$\leftrightarrow$2. Scale by 1. Eliminate all others.
        Record \code{pivot\_cols = \{0,1\}}, \code{cur\_row = 2}.
  \item \textbf{col=2,3:} All entries at/below cur\_row=2 are zero.
        \code{piv = -1} $\Rightarrow$ \code{continue}.
\end{itemize}
Result: rank 2, free columns $\{2,3\}$.
\end{notebox}

\noindent
\textbf{Why Gauss-Jordan (not Gaussian)?}  Plain Gaussian elimination gives
upper triangular $U$.  Gauss-Jordan also eliminates \emph{above} the pivot
(step 4 iterates over \emph{all} rows), giving reduced row echelon form where
each pivot column has exactly one non-zero entry (the leading 1).  This makes
reading off the null-space vectors trivial.

% ─────────────────────────────────────────────────────────────────────────────
\section{LUP Decomposition}
% ─────────────────────────────────────────────────────────────────────────────

\begin{lstlisting}
tuple<Matrix,Matrix,Matrix> Matrix::LUP() const {
    int nn = m;
    Matrix L = eye(nn), U = *this, P = eye(nn);

    for (int k = 0; k < nn; k++) {
        // (1) Partial pivot: find row with largest |entry| in column k
        int piv = k; double best = abs(U.A[k][k]);
        for (int i = k+1; i < nn; i++)
            if (abs(U.A[i][k]) > best) { best = abs(U.A[i][k]); piv = i; }

        if (piv != k) {
            swap(U.A[k], U.A[piv]);    // swap in U
            swap(P.A[k], P.A[piv]);    // mirror swap in P
            // Swap ONLY the already-filled part of L (columns 0..k-1)
            for (int j = 0; j < k; j++) swap(L.A[k][j], L.A[piv][j]);
        }
        if (abs(U.A[k][k]) < EPS) continue;  // singular column

        // (2) Compute multipliers and eliminate below the diagonal
        for (int i = k+1; i < nn; i++) {
            L.A[i][k] = U.A[i][k] / U.A[k][k];   // store multiplier in L
            for (int j = k; j < nn; j++)
                U.A[i][j] -= L.A[i][k] * U.A[k][j];
        }
    }
    return {L, U, P};
}
\end{lstlisting}

\begin{flowbox}{How $PA = LU$ is built}
\begin{itemize}[noitemsep]
  \item $U$ starts as a copy of $A$.  At step $k$, rows $k+1,\ldots,n$ below
        the diagonal in column $k$ are zeroed out.
  \item $L$ starts as $I$.  The multiplier $\ell_{ik} = U_{ik}/U_{kk}$ that
        was used to zero $U_{ik}$ is stored in $L_{ik}$.
  \item $P$ starts as $I$.  Every row swap is mirrored in $P$.  After the
        loop, $P$ encodes the row permutation so $PA = LU$ holds exactly.
  \item The partial swap of $L$ (\code{for (int j = 0; j < k; j++)}) is
        necessary because the previously computed multipliers in $L$ must
        also be permuted consistently.
\end{itemize}
\end{flowbox}

\noindent
\textbf{LUP vs RREF:}  LUP is used only for square matrices to compute
\fn{det} and display $L, U, P$.  Everything else (subspaces, rank, inverse)
uses \fn{rref} because Gauss-Jordan handles non-square matrices naturally.

% ─────────────────────────────────────────────────────────────────────────────
\section{Determinant}
% ─────────────────────────────────────────────────────────────────────────────

\begin{lstlisting}
double Matrix::det() const {
    Matrix U = *this;    // work on a copy
    int swaps = 0;
    for (int k = 0; k < m; k++) {
        int piv = k; double best = abs(U.A[k][k]);
        for (int i = k+1; i < m; i++)
            if (abs(U.A[i][k]) > best) { best = abs(U.A[i][k]); piv = i; }
        if (piv != k) { swap(U.A[k], U.A[piv]); swaps++; }
        if (abs(U.A[k][k]) < EPS) return 0.0;   // singular -> det = 0

        for (int i = k+1; i < m; i++) {
            double f = U.A[i][k] / U.A[k][k];
            for (int j = k; j < m; j++) U.A[i][j] -= f * U.A[k][j];
        }
    }
    double d = (swaps % 2 == 0) ? 1.0 : -1.0;
    for (int i = 0; i < m; i++) d *= U.A[i][i];
    return d;
}
\end{lstlisting}

This is forward-only Gaussian elimination (not full RREF), just enough to get
the upper-triangular form whose diagonal product gives $\det(A)$.
\textbf{Why separate from \fn{LUP}?}  LUP also tracks $L$ and $P$, costing
extra memory.  \fn{det} needs only $U$ and a swap count.

The sign rule: each row swap changes the sign of the determinant, so
$\det(A) = (-1)^s \prod_i U_{ii}$ where $s$ is the total number of swaps.

% ─────────────────────────────────────────────────────────────────────────────
\section{Matrix Inverse via Augmented Gauss-Jordan}
% ─────────────────────────────────────────────────────────────────────────────

\begin{lstlisting}
Matrix Matrix::inv() const {
    int nn = m;
    Matrix aug(nn, 2*nn, 0.0);    // build [A | I]
    for (int i = 0; i < nn; i++) {
        for (int j = 0; j < nn; j++) aug.A[i][j] = A[i][j];
        aug.A[i][nn+i] = 1.0;     // right half = identity
    }
    // Apply Gauss-Jordan to the full 2n-column augmented matrix
    int cur = 0;
    for (int col = 0; col < nn && cur < nn; col++) {
        int piv = -1; double best = EPS;
        for (int r = cur; r < nn; r++)
            if (abs(aug.A[r][col]) > best) { ... piv = r; }
        swap(aug.A[cur], aug.A[piv]);
        // scale and eliminate just as in rref(), but on all 2n columns
        ...
    }
    // Extract right half: that is now A^{-1}
    for (int i = 0; i < nn; i++)
        for (int j = 0; j < nn; j++)
            inv_mat.A[i][j] = aug.A[i][nn+j];
    return inv_mat;
}
\end{lstlisting}

The augmented matrix $[A \mid I]$ is passed through Gauss-Jordan.  Every row
operation applied to $A$ is simultaneously applied to $I$.  When $A$ becomes
$I$ (left half), the right half has become $A^{-1}$.  Every elimination step
operates on \emph{all $2n$ columns} of \code{aug}, which is why the inner loop
uses \code{j < 2*nn}.

% ─────────────────────────────────────────────────────────────────────────────
\section{QR Decomposition (Modified Gram-Schmidt)}
% ─────────────────────────────────────────────────────────────────────────────

\begin{lstlisting}
pair<Matrix,Matrix> Matrix::QR() const {
    Matrix Q = zeros(m, n);   // Q has orthonormal columns
    Matrix R = zeros(n, n);   // R is upper triangular
    Matrix U = *this;         // working copy of A; columns get modified

    for (int j = 0; j < n; j++) {
        // (1) Norm of current j-th column of U
        double norm_sq = 0.0;
        for (int i = 0; i < m; i++) norm_sq += U.A[i][j] * U.A[i][j];
        double norm = sqrt(norm_sq);
        if (norm < EPS) continue;   // linearly dependent column: skip

        R.A[j][j] = norm;
        for (int i = 0; i < m; i++) Q.A[i][j] = U.A[i][j] / norm;

        // (2) KEY: project out Q[:,j] from ALL remaining columns of U immediately
        for (int k = j+1; k < n; k++) {
            double dot = 0.0;
            for (int i = 0; i < m; i++) dot += Q.A[i][j] * U.A[i][k];
            R.A[j][k] = dot;
            for (int i = 0; i < m; i++) U.A[i][k] -= dot * Q.A[i][j];
        }
    }
    return {Q, R};
}
\end{lstlisting}

\begin{notebox}{Classical vs Modified Gram-Schmidt --- the key difference}
\textbf{Classical GS:}
For each column $j$, subtract projections onto $q_1, q_2, \ldots, q_{j-1}$
(the \emph{original} unnormalized vectors from the previous round):
\[
u_j = a_j - \sum_{i<j} \frac{a_j \cdot q_i}{\|q_i\|^2} q_i
\]
This reads from \code{A[i][j]} (unchanged) and subtracts at the end.

\textbf{Modified GS (what we use):}
After computing $q_j$, \emph{immediately} modify all remaining columns of $U$
by subtracting their projection onto $q_j$.  When we later process column $k$,
it has already had $q_1, \ldots, q_{j}$ projected out.  In exact arithmetic,
both give the same result.  But floating-point errors accumulate differently:
MGS reuses the \emph{updated} $U[:, k]$ (which has less of $q_j$ in it), so
each step starts from a vector that is already more orthogonal.  This is why
MGS is preferred in practice.
\end{notebox}

\noindent
\textbf{How $R$ is filled:}
\begin{itemize}[noitemsep]
  \item $R_{jj} = \|u_j\|$ (the diagonal entry records the norm before normalization).
  \item $R_{jk} = q_j \cdot U[:, k]$ for $k > j$ (the inner product with the
        already-normalized $q_j$, stored before subtracting).
  \item $R_{jk} = 0$ for $k < j$ (never written; \code{zeros(n,n)} default).
\end{itemize}

% ─────────────────────────────────────────────────────────────────────────────
\section{Gram-Schmidt Orthonormalization}
% ─────────────────────────────────────────────────────────────────────────────

\fn{gramSchmidt()} is different from \fn{QR}: it operates on the columns of
\code{*this} (which may already be a basis for a subspace) and returns only
the linearly independent columns, orthonormalized.  It is used by all four
\fn{orth*Space()} methods.

\begin{lstlisting}
Matrix Matrix::gramSchmidt() const {
    vector<vector<double>> basis;   // growing list of orthonormal vectors
    int ncols = 0;
    Matrix Q(m, n, 0.0);

    for (int j = 0; j < n; j++) {
        vector<double> v(m);
        for (int i = 0; i < m; i++) v[i] = A[i][j];  // copy column j

        // Subtract projections onto all previously found orthonormal vectors
        for (auto& u : basis) {
            double dot = 0.0;
            for (int i = 0; i < m; i++) dot += u[i] * v[i];
            for (int i = 0; i < m; i++) v[i] -= dot * u[i];
        }
        double norm = 0.0;
        for (int i = 0; i < m; i++) norm += v[i]*v[i];
        norm = sqrt(norm);
        if (norm < EPS) continue;    // linearly dependent -> discard

        for (int i = 0; i < m; i++) v[i] /= norm;
        basis.push_back(v);
        for (int i = 0; i < m; i++) Q.A[i][ncols] = v[i];
        ncols++;
    }
    // Return a matrix with exactly ncols orthonormal columns
    Matrix result(m, ncols);
    for (int i = 0; i < m; i++)
        for (int j = 0; j < ncols; j++)
            result.A[i][j] = Q.A[i][j];
    return result;
}
\end{lstlisting}

\begin{notebox}{Why keep a separate \fn{gramSchmidt()} when \fn{QR()} exists?}
\fn{QR} always returns $Q$ and $R$ of the same size as the input and is
designed for full-matrix factorization.  \fn{gramSchmidt} is simpler: it
discards dependent columns and returns a compact result (fewer columns than the
input if the subspace has lower dimension).  Used as:
\code{nullSpace().gramSchmidt()} $\to$ orthonormal null-space basis.
\end{notebox}

% ─────────────────────────────────────────────────────────────────────────────
\section{Eigenvalue Algorithms}
% ─────────────────────────────────────────────────────────────────────────────

\subsection{\fn{\_eigenSymmetric()} --- Unshifted QR Iteration}

\begin{lstlisting}
EigenResult Matrix::_eigenSymmetric() const {
    int nn = m;
    Matrix T_ = *this;    // will be driven to diagonal form
    Matrix V  = eye(nn);  // accumulates the product Q_1 Q_2 Q_3 ...

    for (int iter = 0; iter < 12000 * nn; iter++) {
        // Convergence check: off-diagonal Frobenius norm
        double off = 0.0;
        for (int i = 0; i < nn; i++)
            for (int j = 0; j < nn; j++)
                if (i != j) off += T_.A[i][j] * T_.A[i][j];
        if (sqrt(off) < EPS * nn) break;

        auto [Q, R] = T_.QR();    // factor T_ = QR
        T_ = R * Q;               // T_{k+1} = R_k Q_k = Q_k^T T_k Q_k
        V  = V * Q;               // V accumulates: V = Q_1 Q_2 ... Q_k
    }
    vector<double> vals(nn), ivals(nn, 0.0);
    for (int i = 0; i < nn; i++) vals[i] = T_.A[i][i];
    return {vals, ivals, V};
}
\end{lstlisting}

\begin{flowbox}{What the iteration does at each step}
\begin{enumerate}[noitemsep]
  \item Compute $T_k = Q_k R_k$ (QR factorization).
  \item Set $T_{k+1} = R_k Q_k$.  Since $T_k = Q_k R_k \Rightarrow
        R_k = Q_k^T T_k$, we have $T_{k+1} = Q_k^T T_k Q_k$ --- a similarity
        transformation.  Eigenvalues are preserved.
  \item Accumulate $V \leftarrow V Q_k$.  After convergence $V$ contains the
        eigenvectors as its columns.
\end{enumerate}
\textbf{Convergence:} for a symmetric matrix, $T_k$ converges to a diagonal
matrix. The off-diagonal Frobenius norm $\|T_k - \operatorname{diag}(T_k)\|_F$
decreases monotonically. The loop breaks when it falls below $\varepsilon n$.
\textbf{Max iterations:} $12000n$ is generous; for a $3\times3$ matrix this is
36\,000 iterations, but convergence typically happens in under 300.
\end{flowbox}

\noindent
\textbf{Why no Rayleigh shift?}  A common speed-up is to subtract $\mu I$
(where $\mu = T_{k,nn-1,nn-1}$) before the QR step and add it back after.
When $\mu$ is exactly an eigenvalue (e.g., a matrix with repeated eigenvalue 2
at the bottom-right), $T_k - \mu I$ becomes \emph{singular}: one column has
near-zero norm in QR and the corresponding $Q$ column is never computed,
leaving a zero column in $V$.  The unshifted version avoids this completely.

\subsection{\fn{\_eigenGeneral()} --- Faddeev-LeVerrier + Durand-Kerner}

This function handles non-symmetric matrices, which may have complex eigenvalues.
It has two phases.

\subsubsection{Phase 1: Faddeev-LeVerrier --- computing the characteristic polynomial}

\begin{lstlisting}
vector<double> coeff(nn + 1);
coeff[nn] = 1.0;        // leading coefficient is always 1
Matrix B = eye(nn);
for (int k = 1; k <= nn; k++) {
    B = (*this) * B;    // B <- A * B_{k-1}   (a fresh matrix multiply)
    double tr = 0.0;
    for (int i = 0; i < nn; i++) tr += B.A[i][i];   // trace of current B
    coeff[nn - k] = -tr / k;
    if (k < nn)
        for (int i = 0; i < nn; i++) B.A[i][i] += coeff[nn - k];  // B <- B + c*I
}
\end{lstlisting}

After the loop, \code{coeff[0..n]} hold the coefficients of
$p(\lambda) = \lambda^n + c_{n-1}\lambda^{n-1} + \cdots + c_0$.
The recurrence is: $c_{n-k} = -\operatorname{tr}(A \cdot B_{k-1}) / k$ and
$B_k = A B_{k-1} + c_{n-k} I$.  This uses $n$ matrix multiplications.

\subsubsection{Phase 2: Durand-Kerner --- finding all $n$ roots simultaneously}

\begin{lstlisting}
// Lambda for polynomial evaluation using Horner's method
auto poly_eval = [&](C z) -> C {
    C r(coeff[nn]);                                    // start with leading 1
    for (int k = nn - 1; k >= 0; k--) r = r * z + C(coeff[k]);
    return r;                                          // p(z)
};

// Initial positions: n points on a circle of radius r0 (Cauchy bound)
double cauchy = 0.0;
for (int k = 0; k < nn; k++)
    cauchy = max(cauchy, abs(coeff[k] / coeff[nn]));
double r0 = 1.0 + cauchy;
vector<C> roots(nn);
for (int k = 0; k < nn; k++) {
    double ang = 2.0 * M_PI * k / nn + 0.17;
    roots[k] = C(r0 * cos(ang), r0 * sin(ang));
}

for (int iter = 0; iter < 1000 * nn; iter++) {
    double max_chg = 0.0;
    vector<C> nr(nn);
    for (int i = 0; i < nn; i++) {
        // Denominator: product of (z_i - z_j) for all j != i
        C denom(1.0);
        for (int j = 0; j < nn; j++)
            if (j != i) denom *= (roots[i] - roots[j]);
        if (abs(denom) < 1e-15) { nr[i] = roots[i]; continue; }  // guard
        C corr = poly_eval(roots[i]) / denom;
        nr[i] = roots[i] - corr;        // Weierstrass correction
        max_chg = max(max_chg, abs(corr));
    }
    roots = nr;
    if (max_chg < 1e-12) break;
}
\end{lstlisting}

\begin{flowbox}{How Durand-Kerner works step by step}
\begin{itemize}[noitemsep]
  \item \textbf{Polynomial evaluation:} \fn{poly\_eval} uses Horner's method
        $(\ldots((c_n z + c_{n-1})z + c_{n-2})\ldots)z + c_0)$ which is
        numerically stable and requires only $n$ multiplications.
  \item \textbf{Initial guess radius} \code{r0}: the Cauchy bound
        $1 + \max_k |c_k / c_n|$ guarantees all roots lie \emph{inside} the
        circle $|z| \leq r_0$.  Starting outside ensures no roots are missed.
  \item \textbf{Initial angle offset} \code{0.17}: placing roots at exact
        equally-spaced angles would cause symmetry that makes the denominator
        product zero; the \code{0.17} radian offset breaks this.
  \item \textbf{The correction} $p(z_i) / \prod_{j\neq i}(z_i - z_j)$ is the
        Newton step for the deflated polynomial $p(z) / \prod_{j \neq i}
        (z - z_j)$, whose unique root near $z_i$ is the actual eigenvalue.
        All $n$ corrections are computed \emph{before} any update (stored in
        \code{nr}), making it a Jacobi-style simultaneous update.
  \item \textbf{Guard} \code{if (abs(denom) < 1e-15)}: if two initial guesses
        are identical (very rare), the denominator would be zero.  We skip the
        update and let the iteration sort them out.
\end{itemize}
\end{flowbox}

\subsubsection{Phase 3: Classify and clean up eigenvalues}

\begin{lstlisting}
// Classify: imaginary part is "zero" if |Im| < 1e-4 * max(1, |Re|)
for (int i = 0; i < nn; i++) {
    double re = roots[i].real(), im = roots[i].imag();
    bool real_eig = abs(im) < max(1e-6, 1e-4 * abs(re));
    real_vals[i] = (abs(re) < 1e-9) ? 0.0 : re;
    imag_vals[i] = real_eig ? 0.0 : im;
}
\end{lstlisting}

For repeated real eigenvalues (e.g., $\lambda = 2$ with multiplicity 3),
Durand-Kerner may not fully converge and leaves tiny imaginary parts $\sim
10^{-5}$.  The \emph{relative} threshold \code{1e-4 * |re|} catches these
(for $\lambda \approx 2$, the threshold is $2 \times 10^{-4}$, larger than
$10^{-5}$) without misclassifying genuinely complex eigenvalues (where
$|\text{Im}| \sim |\text{Re}|$).

\subsubsection{Phase 4: Eigenvectors via \fn{nullSpace}}

\begin{lstlisting}
Matrix V_ = eye(nn);  // fallback if all eigenvalues are complex
for (int k = 0; k < nn; k++) {
    if (abs(imag_vals[k]) > 1e-7) continue;  // skip complex eigenvalues
    double lam = real_vals[k];
    Matrix A_lam = *this;
    for (int i = 0; i < nn; i++) A_lam.A[i][i] -= lam;  // A - lambda*I
    Matrix ns = A_lam.nullSpace();   // RREF -> free variable vectors
    if (ns.n > 0)
        for (int i = 0; i < nn; i++) V_.A[i][k] = ns.A[i][0];
}
\end{lstlisting}

For each real eigenvalue $\lambda$, form $(A - \lambda I)$ by subtracting
$\lambda$ from the diagonal, then call \fn{nullSpace()} which runs RREF to
find the null-space basis.  The first basis vector is stored as the $k$-th
column of $V\_$.

\subsection{\fn{eigen()} --- the dispatcher}

\begin{lstlisting}
EigenResult Matrix::eigen() const {
    if (!isSquare()) throw runtime_error("eigen: not square");
    if (isSymmetric()) return _eigenSymmetric();
    return _eigenGeneral();
}
\end{lstlisting}

\fn{isSymmetric()} checks $|A_{ij} - A_{ji}| < \varepsilon$ for all entries.
For symmetric matrices the unshifted QR iteration is used (guaranteed real
eigenvalues, correct eigenvectors).  For non-symmetric matrices, Faddeev-LeVerrier
+ Durand-Kerner handles the general case including complex eigenvalues.

% ─────────────────────────────────────────────────────────────────────────────
\section{SVD}
% ─────────────────────────────────────────────────────────────────────────────

\begin{lstlisting}
tuple<Matrix,Matrix,Matrix> Matrix::SVD() const {
    Matrix AtA = T() * (*this);        // A^T A  (n x n, symmetric PSD)
    auto er = AtA._eigenSymmetric();   // eigenvalues in er.real_vals, eigenvectors in er.vecs
    int rr = min(m, n);

    // Sort eigenpairs by eigenvalue descending (largest sigma first)
    vector<pair<double,int>> sp;
    for (int i = 0; i < n; i++) sp.push_back({er.real_vals[i], i});
    sort(sp.begin(), sp.end(), [](auto& a, auto& b){ return a.first > b.first; });

    Matrix V(n, n, 0.0);
    vector<double> sigmas(n);
    for (int j = 0; j < n; j++) {
        int idx = sp[j].second;
        for (int i = 0; i < n; i++) V.A[i][j] = er.vecs.A[i][idx]; // reorder V cols
        sigmas[j] = sqrt(max(0.0, sp[j].first));   // sigma_j = sqrt(lambda_j)
    }

    // Build Sigma (m x n diagonal)
    Matrix S = zeros(m, n);
    for (int i = 0; i < rr; i++) S.A[i][i] = sigmas[i];

    // Compute U columns: u_j = A v_j / sigma_j
    Matrix U(m, m, 0.0);
    int rank_A = 0;
    for (int j = 0; j < rr; j++) {
        if (sigmas[j] > EPS) {
            for (int i = 0; i < m; i++) {
                double s = 0.0;
                for (int k = 0; k < n; k++) s += A[i][k] * V.A[k][j];
                U.A[i][rank_A] = s / sigmas[j];   // u = Av / sigma
            }
            rank_A++;
        }
    }

    // Complete U: add orthonormal basis for null(A^T) using Gram-Schmidt
    // against the already-computed u columns
    vector<vector<double>> u_basis;
    for (int j = 0; j < rank_A; j++) { /* collect existing U columns */ }
    for (int k = 0; k < m && cur_col < m; k++) {
        vector<double> e(m, 0.0); e[k] = 1.0;  // try each standard basis vector
        // Orthogonalize e against all vectors in u_basis
        for (auto& u : u_basis) {
            double dot = ...; for (...) e[i] -= dot * u[i];
        }
        double norm = ...; if (norm < EPS) continue;
        // Accept, normalize, add to u_basis and U
    }
    return {U, S, V};
}
\end{lstlisting}

\begin{flowbox}{SVD computation chain}
\begin{enumerate}[noitemsep]
  \item Form $B = A^TA$ (symmetric, positive semi-definite).
  \item Call \fn{\_eigenSymmetric()} on $B$.  Returns eigenvalues
        $\lambda_1 \geq \cdots \geq \lambda_n \geq 0$ and eigenvectors
        (columns of $V$).
  \item $\sigma_j = \sqrt{\lambda_j}$.  Build $\Sigma$ (diagonal, $m \times n$).
  \item For each $\sigma_j > \varepsilon$: $u_j = A v_j / \sigma_j$.  This is
        the relationship $Av = \sigma u$.
  \item Complete $U$ to $m \times m$ orthogonal by Gram-Schmidt:  try each
        standard basis vector $e_k$, orthogonalize against the already-computed
        $u_j$'s; keep it if the norm after orthogonalization is $> \varepsilon$.
\end{enumerate}
\end{flowbox}

% ─────────────────────────────────────────────────────────────────────────────
\section{Pseudoinverse}
% ─────────────────────────────────────────────────────────────────────────────

\begin{lstlisting}
Matrix Matrix::pinv() const {
    auto [U, S, V] = SVD();
    Matrix Sp = zeros(n, m);      // Sigma^+ is n x m (transposed shape)
    for (int i = 0; i < min(m,n); i++)
        if (S.A[i][i] > EPS) Sp.A[i][i] = 1.0 / S.A[i][i];  // invert nonzero sigmas
    return V * Sp * U.T();        // A^+ = V Sigma^+ U^T
}
\end{lstlisting}

\textbf{Shape of $\Sigma^+$:}  $\Sigma$ is $m \times n$; its pseudoinverse
$\Sigma^+$ is $n \times m$ (transposed shape).  If $\Sigma_{ii} = \sigma_i$,
then $\Sigma^+_{ii} = 1/\sigma_i$ (for nonzero $\sigma_i$).  This is then
used in $V(n \times n) \cdot \Sigma^+(n \times m) \cdot U^T(m \times m) =
A^+(n \times m)$.

% ─────────────────────────────────────────────────────────────────────────────
\section{Fundamental Subspaces}
% ─────────────────────────────────────────────────────────────────────────────

\subsection{\fn{nullSpace()} --- from free variables in RREF}

\begin{lstlisting}
Matrix Matrix::nullSpace() const {
    auto [R, pivot_cols, rank_] = rref();
    int num_free = n - rank_;
    if (num_free == 0) return Matrix(n, 0);   // trivial null space

    // Mark which columns are pivot vs free
    vector<bool> is_pivot(n, false);
    for (int c : pivot_cols) is_pivot[c] = true;
    vector<int> free_cols;
    for (int j = 0; j < n; j++) if (!is_pivot[j]) free_cols.push_back(j);

    Matrix N(n, num_free, 0.0);
    for (int fc = 0; fc < num_free; fc++) {
        int fc_col = free_cols[fc];
        N.A[fc_col][fc] = 1.0;   // set free variable = 1
        for (int pr = 0; pr < rank_; pr++)
            N.A[pivot_cols[pr]][fc] = -R.A[pr][fc_col];  // pivot = -(RREF entry)
    }
    return N;
}
\end{lstlisting}

\begin{notebox}{How null-space vectors are read from RREF}
From the RREF $R$, each free column $f$ gives one null-space vector.
The equation $Rx = 0$ means: for each pivot row $p_k$:
$x_{p_k} + \sum_{\text{free }f} R_{p_k,f} x_f = 0$.
Setting $x_f = 1$ (for this specific free column) and all other free
variables to 0 gives $x_{p_k} = -R_{p_k, f}$.

Example for RREF $\begin{pmatrix}1&0&-1&-2\\0&1&2&3\\0&0&0&0\end{pmatrix}$:
\begin{itemize}[noitemsep]
  \item Pivot cols: $\{0,1\}$.  Free cols: $\{2,3\}$.
  \item Free col 2 ($x_2=1, x_3=0$): $x_0=-(-1)=1$, $x_1=-(2)=-2$.
        Vector: $[1,-2,1,0]^T$.
  \item Free col 3 ($x_2=0, x_3=1$): $x_0=-(-2)=2$, $x_1=-(3)=-3$.
        Vector: $[2,-3,0,1]^T$.
\end{itemize}
\end{notebox}

\subsection{\fn{columnSpace()} --- original pivot columns}

\begin{lstlisting}
Matrix Matrix::columnSpace() const {
    auto [R, pivot_cols, rank_] = rref();
    Matrix C(m, rank_);
    for (int j = 0; j < rank_; j++)
        for (int i = 0; i < m; i++)
            C.A[i][j] = A[i][pivot_cols[j]];   // original A, not R
    return C;
}
\end{lstlisting}

\textbf{Why use original $A$, not RREF $R$?}  The pivot columns of $R$ always
have $1$ in one row and $0$ elsewhere (standard unit vectors).  That's not a
useful basis for the column space of the original $A$.  The pivot columns of
the \emph{original} $A$ span the column space and represent geometrically
meaningful directions.

\subsection{\fn{rowSpace()} --- pivot rows of RREF as columns}

\begin{lstlisting}
Matrix Matrix::rowSpace() const {
    auto [R, pivot_cols, rank_] = rref();
    Matrix RS(n, rank_);
    for (int i = 0; i < rank_; i++)
        for (int j = 0; j < n; j++)
            RS.A[j][i] = R.A[i][j];   // transpose: row i of R -> column i of RS
    return RS;
}
\end{lstlisting}

The nonzero rows of RREF are the canonical basis for the row space.  They are
transposed (stored as \emph{columns} of the returned matrix) so the
conventions match: all subspace bases are returned as matrices whose
\emph{columns} span the subspace.

\subsection{\fn{leftNullSpace()} --- one-liner}

\begin{lstlisting}
Matrix Matrix::leftNullSpace() const { return T().nullSpace(); }
\end{lstlisting}

$\mathcal{N}(A^T)$ is the null space of $A^T$.  Transpose $A$, call
\fn{nullSpace()}, done.

\subsection{Orthonormal versions}

\begin{lstlisting}
Matrix Matrix::orthNullSpace() const {
    auto N = nullSpace();
    return (N.n == 0) ? N : N.gramSchmidt();
}
\end{lstlisting}

The pattern is the same for all four: compute the (possibly non-orthogonal)
basis, then call \fn{gramSchmidt()} on it if it is non-trivial.  The
\code{N.n == 0} guard handles the trivial (zero) subspace.

% ─────────────────────────────────────────────────────────────────────────────
\section{Projection Matrices}
% ─────────────────────────────────────────────────────────────────────────────

\begin{lstlisting}
Matrix Matrix::projOntoColSpace()  const { return (*this) * pinv(); }
Matrix Matrix::projOntoRowSpace()  const { return T().projOntoColSpace().T(); }
Matrix Matrix::projOntoNullSpace() const { return eye(m) - projOntoColSpace(); }
\end{lstlisting}

\begin{flowbox}{Why these formulas work}
\textbf{Onto column space:} $P_C = A A^+$.  For any $b$: $P_C b = A A^+ b = A
\hat{x}$ where $\hat{x} = A^+ b$ is the least-squares solution.  $A\hat{x}$
is the closest point in $\Col(A)$ to $b$.

\textbf{Onto null space:} $P_N = I - P_C$.  Any vector decomposes as
$b = P_C b + P_N b$, the column-space component plus the null-space component.

\textbf{Onto row space:} The row space of $A$ is the column space of $A^T$.
So $P_R = (A^T)(A^T)^+ = A^T (A^{T+})$. But $(A^T)^+ = (A^+)^T$ (a property
of the pseudoinverse), so $P_R = A^T (A^+)^T = (A^+ A)$.  In code:
\code{T().projOntoColSpace().T()} computes $(A^T (A^T)^+)^T = (A^+ A) = P_R$.
\end{flowbox}

% ─────────────────────────────────────────────────────────────────────────────
\section{Solve \texttt{Ax = b}}
% ─────────────────────────────────────────────────────────────────────────────

\begin{lstlisting}
Matrix Matrix::solve(const Matrix& b) const {
    if (b.n != 1 || b.m != m) throw runtime_error("solve: dimension mismatch");
    Matrix x = pinv() * b;
    // Snap near-zero solution components to exact zero
    for (int i = 0; i < x.m; i++)
        for (int j = 0; j < x.n; j++)
            if (abs(x.A[i][j]) < EPS) x.A[i][j] = 0.0;
    return x;
}
\end{lstlisting}

The minimum-norm least-squares solution is $\hat{x} = A^+ b$.  This works
in all cases:
\begin{itemize}[noitemsep]
  \item $A$ full rank, square: $A^+ = A^{-1}$, so $\hat{x} = A^{-1}b$ (exact solution).
  \item $b \in \Col(A)$, $A$ rank-deficient: $\hat{x}$ is the minimum-norm
        exact solution.
  \item $b \notin \Col(A)$: $\hat{x}$ minimises $\|Ax - b\|_2$.
\end{itemize}

% ─────────────────────────────────────────────────────────────────────────────
\section{Jordan Form}
% ─────────────────────────────────────────────────────────────────────────────

\begin{lstlisting}
JordanInfo Matrix::jordanForm() const {
    auto er = eigen();
    vector<bool> seen(nn, false);

    for (int i = 0; i < nn; i++) {
        if (seen[i]) continue;
        seen[i] = true;
        int cnt = 1;
        double scale_i = max(1.0, abs(er.real_vals[i]));
        // Cluster nearby eigenvalues into one group
        for (int j = i+1; j < nn; j++) {
            if (!seen[j]
                && abs(er.real_vals[j] - er.real_vals[i]) < 1e-3 * scale_i
                && abs(er.imag_vals[j] - er.imag_vals[i]) < 1e-3 * scale_i) {
                seen[j] = true; cnt++;
            }
        }
        alg_m.push_back(cnt);   // algebraic multiplicity

        // Geometric multiplicity: dim null(A - lambda*I) = n - rank(A - lambda*I)
        if (abs(er.imag_vals[i]) < EPS) {
            Matrix AlamI = *this;
            for (int k = 0; k < nn; k++) AlamI.A[k][k] -= er.real_vals[i];
            int gm = max(1, nn - AlamI.rank());   // rank() calls rref() internally
            geo_m.push_back(gm);
        } else {
            geo_m.push_back(1);  // complex: assume one block
        }
    }

    // Build Jordan form matrix J block by block
    int pos = 0;
    for (int k = 0; k < uniq_r.size(); k++) {
        int am = alg_m[k], nb = geo_m[k];   // nb = num blocks = geo mult
        int base = am / nb, extra = am % nb;
        for (int b = 0; b < nb; b++) {
            int bsz = base + (b < extra ? 1 : 0);  // distribute am into nb blocks
            for (int r = 0; r < bsz; r++) {
                J.A[pos+r][pos+r] = lam;               // diagonal = lambda
                if (r < bsz-1) J.A[pos+r][pos+r+1] = 1.0; // superdiagonal = 1
            }
            pos += bsz;
        }
    }
}
\end{lstlisting}

\begin{notebox}{How block sizes are determined}
Given algebraic multiplicity $m_a$ and geometric multiplicity (number of blocks)
$m_g$:
\begin{itemize}[noitemsep]
  \item $m_a$ divided equally into $m_g$ blocks: base size $= \lfloor m_a /
        m_g \rfloor$, first \code{extra = m\_a \% m\_g} blocks get $+1$.
  \item Example: $m_a = 5$, $m_g = 2 \Rightarrow$ blocks of size 3 and 2.
  \item Example: $m_a = 3$, $m_g = 1 \Rightarrow$ one block of size 3 (single
        Jordan chain, the defective case).
\end{itemize}
\end{notebox}

% ─────────────────────────────────────────────────────────────────────────────
\section{Change of Basis and Similar Matrices}
% ─────────────────────────────────────────────────────────────────────────────

\begin{lstlisting}
Matrix Matrix::changeOfBasis(const Matrix& P) const {
    return P.inv() * (*this) * P;   // P^{-1} A P
}

bool Matrix::isSimilarTo(const Matrix& B) const {
    auto sort_ev = [](vector<double> r, vector<double> im) {
        vector<pair<double,double>> v;
        for (int i = 0; i < r.size(); i++) v.push_back({r[i], im[i]});
        sort(v.begin(), v.end());
        return v;
    };
    auto ea = eigen(), eb = B.eigen();
    if (ea.real_vals.size() != eb.real_vals.size()) return false;
    auto sa = sort_ev(ea.real_vals, ea.imag_vals);
    auto sb = sort_ev(eb.real_vals, eb.imag_vals);
    for (int i = 0; i < sa.size(); i++) {
        if (abs(sa[i].first  - sb[i].first)  > 1e-5) return false;
        if (abs(sa[i].second - sb[i].second) > 1e-5) return false;
    }
    return abs(det() - B.det()) < 1e-5;
}
\end{lstlisting}

\fn{isSimilarTo} sorts both eigenvalue lists (as complex pairs by real part)
and compares them element-wise.  The determinant is also compared as an extra
check.  This is a necessary condition, not sufficient in general, but works
for the matrices typical in a linear algebra course.

% ─────────────────────────────────────────────────────────────────────────────
\section{Function Call Graph}
% ─────────────────────────────────────────────────────────────────────────────

The diagram below shows which functions call which others.
Arrows mean ``calls''.

\begin{center}
\small
\begin{tabular}{ll}
\toprule
Caller & Callees \\
\midrule
\fn{solve}        & \fn{pinv} \\
\fn{pinv}         & \fn{SVD} \\
\fn{SVD}          & \fn{T}, \fn{operator*}, \fn{\_eigenSymmetric}, \fn{zeros}, \fn{gramSchmidt*} \\
\fn{\_eigenSymmetric} & \fn{QR}, \fn{operator*} \\
\fn{\_eigenGeneral}  & \fn{operator*} (F-L), \fn{nullSpace} (eigenvectors) \\
\fn{eigen}        & \fn{isSymmetric}, \fn{\_eigenSymmetric} or \fn{\_eigenGeneral} \\
\fn{jordanForm}   & \fn{eigen}, \fn{rank} \\
\fn{rank}         & \fn{rref} \\
\fn{nullSpace}    & \fn{rref} \\
\fn{columnSpace}  & \fn{rref} \\
\fn{rowSpace}     & \fn{rref} \\
\fn{leftNullSpace} & \fn{T}, \fn{nullSpace} \\
\fn{orth*Space}   & \fn{*Space}, \fn{gramSchmidt} \\
\fn{projOntoColSpace} & \fn{operator*}, \fn{pinv} \\
\fn{projOntoRowSpace} & \fn{T}, \fn{projOntoColSpace} \\
\fn{projOntoNullSpace} & \fn{eye}, \fn{projOntoColSpace} \\
\fn{det}          & (self-contained) \\
\fn{inv}          & \fn{isSingular} ($\to$ \fn{det}) \\
\fn{QR}           & \fn{zeros}, \fn{operator*} (implicitly) \\
\fn{LUP}          & \fn{eye}, \fn{isSquare} \\
\fn{isPositiveDefinite} & \fn{isSymmetric}, \fn{\_eigenSymmetric} \\
\fn{changeOfBasis} & \fn{inv}, \fn{operator*} \\
\fn{isSimilarTo}  & \fn{eigen}, \fn{det} \\
\bottomrule
\end{tabular}
\end{center}

\begin{notebox}{Central dependency: \fn{rref()}}
Almost every subspace function calls \fn{rref()}.  Since \fn{rref()} makes a
\emph{copy} (\code{Matrix R = *this}) it never mutates the original object,
so repeated calls are safe (though not cached).
\end{notebox}

% ─────────────────────────────────────────────────────────────────────────────
\section{main.cpp --- Program Flow}
% ─────────────────────────────────────────────────────────────────────────────

\subsection{Input phase}

\begin{lstlisting}
int m, n;
cout << "Enter matrix dimensions (rows cols): ";
cin >> m >> n;
Matrix A = Matrix::fromInput(m, n);
\end{lstlisting}

\fn{fromInput} reads each entry from \code{cin}.  After this, \code{A} is the
only data structure; everything else is derived from it.

\subsection{Output sections --- always computed}

Regardless of shape ($m \times n$, any rank):
\begin{enumerate}[noitemsep]
  \item \textbf{RREF + pivots:} \code{auto rres = A.rref();}
  \item \textbf{Four fundamental subspaces} + \textbf{orthonormal bases}.
  \item \textbf{QR} (only if $m \geq n$): \code{auto [Q,R] = A.QR();}
  \item \textbf{SVD}: \code{auto [U,S,V] = A.SVD();}
  \item \textbf{Projection matrices:} \fn{projOntoColSpace}, \fn{projOntoNullSpace}.
\end{enumerate}

\subsection{Square-matrix section}

Gated on \code{if (A.isSquare())}:
\begin{enumerate}[noitemsep]
  \item \textbf{LUP}: \code{auto [L,U,P] = A.LUP();}
  \item \textbf{Properties}: symmetric, orthogonal, singular, PD, PSD.
  \item \textbf{Determinant}: \code{A.det();}
  \item \textbf{Inverse or pseudoinverse}: \code{A.inv()} if not singular,
        else \code{A.pinv()}.
  \item \textbf{Eigenvalues/vectors}: \code{auto er = A.eigen();}
  \item \textbf{Jordan form}: \code{auto ji = A.jordanForm();}
  \item \textbf{Diagonalization demo}: if all $m_a = m_g$ and $n \leq 5$,
        compute $P^{-1}AP$ with $P =$ eigenvector matrix.
\end{enumerate}

For \emph{non-square} matrices, the pseudoinverse and row-space projection are
additionally shown.

\subsection{Solve phase}

\begin{lstlisting}
cout << "Enter column vector b (" << m << " entries): ";
Matrix b_vec(m, 1);
for (int i = 0; i < m; i++) cin >> b_vec.A[i][0];

Matrix x_sol = A.solve(b_vec);
// Compute and print residual
Matrix residual = A * x_sol - b_vec;
double res = 0.0;
for (int i = 0; i < m; i++) res += residual.A[i][0] * residual.A[i][0];
res = sqrt(res);
\end{lstlisting}

After displaying all structural properties, the program asks for $b$ and
calls \fn{solve}.  The residual $\|Ax - b\|_2$ confirms whether an exact
solution was found or the result is only least-squares.

% ─────────────────────────────────────────────────────────────────────────────
\section{Numerical Constants and Tolerances Summary}
% ─────────────────────────────────────────────────────────────────────────────

\begin{center}
\small
\renewcommand{\arraystretch}{1.4}
\begin{tabular}{lll}
\toprule
Constant & Value & Used for \\
\midrule
\code{EPS} & $10^{-9}$ & Zero threshold everywhere (pivots, norms, det) \\
\code{1e-6} & $10^{-6}$ & \fn{approxEqual} (orthogonality check) \\
\code{1e-5} & $10^{-5}$ & \fn{isSimilarTo} eigenvalue comparison \\
\code{1e-3 * scale\_i} & relative & Jordan form eigenvalue clustering \\
\code{1e-4 * |re|} & relative & Classifying complex vs real eigenvalue \\
\code{1e-12} & $10^{-12}$ & Durand-Kerner convergence (max correction) \\
\code{1e-15} & $10^{-15}$ & Durand-Kerner denominator guard \\
\code{12000 * n} & --- & Max QR iterations for symmetric eigenvalues \\
\code{1000 * n} & --- & Max Durand-Kerner iterations \\
\bottomrule
\end{tabular}
\end{center}

% ─────────────────────────────────────────────────────────────────────────────
\section{Known Limitations in the Code}
% ─────────────────────────────────────────────────────────────────────────────

\begin{enumerate}[noitemsep]
  \item \textbf{No caching:} \fn{rref()} is called independently by
        \fn{nullSpace}, \fn{columnSpace}, \fn{rowSpace}, and \fn{rank}.
        For large matrices this is wasteful; a memoised RREF result would
        help.
  \item \textbf{No complex eigenvectors:} \fn{\_eigenGeneral} skips
        eigenvectors for complex eigenvalues (\code{if (abs(imag) > 1e-7)
        continue}).  A full treatment would use $\mathbb{C}$ arithmetic.
  \item \textbf{Rank of $(A - \lambda I)$ via RREF:}  For nearly-repeated
        eigenvalues, the floating-point $\lambda$ may not be accurate enough
        for the RREF rank to be correct.  The \code{max(1, nn - rank)} guard
        prevents geo\_mult from being zero, but the Jordan block structure may
        be wrong if eigenvalues are very close.
  \item \textbf{$O(n^4)$ worst case:}  \fn{jordanForm} calls \fn{rank}
        (which calls \fn{rref}, $O(n^3)$) once per unique eigenvalue.  For
        a matrix with $n$ distinct eigenvalues, total cost is $O(n^4)$.
  \item \textbf{QR on square submatrix:}  \fn{\_eigenSymmetric} passes the
        full $n \times n$ matrix to \fn{QR} at every iteration.
        A proper implementation would reduce to tridiagonal form once and
        then do QR steps on a tridiagonal, saving a factor of $n$.
\end{enumerate}

\end{document}
